\clearpage
\part{Unresolved questions and unchecked source assertions}
\section{Current gates and approaches still in progress}
No mathematical assertion in this part is promoted to a theorem by
its inclusion. The executable question register is \texttt{QUESTIONS.tsv};
the dated approach record is \texttt{LEDGER.md}. A missing historical
proof package and an open mathematical conjecture are different statuses.

\begin{description}[style=nextline,leftmargin=0pt]
\item[Cote 14, scans 25--38: original quotient identification.]
The generated quotient has full proofs in AQ-14-FORWARD and
AQ-14-CONVERSE. Original Definition1.1, Proposition1.2 and
Corollary2.3 now identify that quotient under the explicitly bracketed
projective-parameter convention. AQ-14-QUASIPROJECTIVE now removes the projective qualification
in characteristic zero for smooth quasiprojective parameters, by an
integral equality of the generated relations. AQ-14-NONPROPER now
proves the cohomological implication in every characteristic.
AQ-14-ALTERATION proves that the relation quotient in characteristic
$p>0$ is $p$-primary torsion, hence proves equality after inverting $p$
and on the prime-to-$p$ torsion in the resulting cycle quotients.
The possible residual $p$-primary relation difference is the remaining
integral question. Supplied scan38 is physical PDFpage38,
handwritten page37; both locators are retained.
\item[Cote 15, scans 100--106: full simple-motive classification.]
Weil-number occurrence is proved. The cited 2026 CM Hodge theorem,
with the inspected Milne implications, discharges the rational Tate and
standard-conjecture inputs for abelian varieties and their products.
AQ-ABELIAN-RECOGNITION proves full faithfulness after scalar extension
and Frobenius-polynomial recognition in the abelian--Artin numerical
category. CC-AG01--03 further classify its idempotent completion over
$\overline{\mathbf Q}_\ell$ by finite Weil-number gradings and
compute its diagonalizable group, including its Artin torsion quotient. CC-XC01--04 now reduce membership of a specified smooth projective
variety to a determinant condition on one projected finite family
of cross-cycle maps. This condition forces semisimplicity and yields
fixed-prime Tate and numerical--homological comparison on the stated
products. Extending the classification to arbitrary varieties still
requires producing those actual cross-correspondences. The
quaternion example AQ-RATIONAL-FIBRE also shows why a rational neutral
fibre functor and arbitrary rational multiplicities cannot be silently
inserted. AQ-WEIL-NOTATION recovers the historical totally real and CM
closures on the earlier original sheets and proves the norm-condition
identification; that historical notation is no longer a gate.
The hypotheses of Milne (1994), \S2, have therefore been
discharged in a specified sector, not deleted globally.
\item[Cote 18, scans 11--13 and 59: period density.]
The overly strong formulation is refuted, and generation for one
elliptic curve has a precise known answer. Non-CM transcendence degree
four and the non-CM algebraicity question for the Hodge-splitting
scalar remain the specified transcendence gates. Group generation
does not imply density of a particular period locus; AQ-COCYCLE
gives the exact distinction.
\item[Cote 134-8, scans 47--54: the linear theorem.]
The original predicates and adjunction directions have been read.
BR-13401--06 and F-SM-17 prove the precise modern comparison and
counit criteria, the equivalence over fields and \(\mathbf Z\) under
the stated hypotheses, and a locally projective ringoid construction.
F-SM-24 extends this to all countable rings and all scalar rings of
cardinality $\aleph_t$ for finite $t$, under the explicit scalar
compatibility, and to arbitrary scalar rings when the ringoid morphism
modules are finitely generated projective. F-SM-27 gives a third range:
arbitrary scalar rings and projective morphism modules if the left
module $D$ has a resolution by finitely generated projectives.
Outside these ranges the remaining test is
\(L\varphi(GP)=0\) on acyclic projective chain objects, equivalently
the explicit Tor condition SM25.2. F-SM-26 shows why finite flat
dimension of $D$ cannot furnish another case: it contradicts the
nonzero full-interval normalization. The intake's
three-sort counterexample remains an uninspected earlier construction;
no claim about its exact compatibility with the original hypotheses
has been promoted to a new proof here.
\item[SA08, cote 136/117--121: strict PD sources and localization.]
AQ-SA08-LIFT refutes a specific degreewise-free lifting input.
The original admissibility conditions and common-source and strict
equalization conjectures have now been retrieved on physical pages
111, 113, 115, 117, 119 and 121, footers one smaller. Their exact
transcription and a proof or counterexample under those conditions
are in progress. The stronger intake all-arrow terminality and SA09
comparison still require their distinct category definitions and proof
context. The lifting counterexample is not a refutation of those
different existence questions.
\item[Cote 123: apposition outside type A.]
The quotient-dimension, rational-singularity and Lie-rank questions
are answered. CC-RK01--05 supplies the remaining classical and
exceptional ranks, subregular uniqueness and the analytic transfer
to unipotent group elements. BR-12305--08 treats all central forms
of the group quotient. CC-AP01--07 completely classifies type A
apposition quadruples, their generators, opposition and outer action.
Explicit classifications in the other simple types remain in progress;
type A already refutes a universal distinguished outer-invariant class.
\item[Weight--monodromy and universal-curve categories.]
The local monodromy operator and its invariant exterior powers must
be retained. BR-WM-02 supplies a finite gate and has passed a
recorded independent whole-proof reading. CC-ST01--04 give the geometric
stack decomposition with its explicit separating-divisor summands.
CC-RES01--02 prove the canonical exit/right-Kan statement and the nodal
lifting-category calculation. Identifying an unseen chosen archive
functor with that canonical functor, or deducing all categorical base
changes from geometric proper base change, requires another argument.
Purity of a proper pushforward
complex is not purity of every ordinary singular-fibre cohomology group.
\item[The analytic method as a transferable theorem.]
The lattice ideals, repeated theta-lift estimates, period nonvanishing,
and tower bounds must be exported with exact hypotheses from the
full argument. BR-DET-01/02 identify the numerical determinant gate.
The footnote on $K$-groups contains no degree-indexed theorem to apply.
F-SM-14--16 instead give exact comparisons and explicit families of
torsion consequences of separately cited integral arithmetic theorems.
\item[Masullo's operator realization.]
The fully normalized numerical shrinking-strip identity is proved in
the Dirichlet case in F-SM-13, including the finite Fourier expansion
of the tests and the archimedean parity term. F-SM-18 proves diagonal summation with a valid normalized indentation,
and F-SM-20 gives the actual compact-invariant compression. F-SM-19
refutes the printed indentation: it vanishes for $\mathbf Q$.
F-SM-21--22 rule out two specific naive replacements, by a nonzero
prime-trace discrepancy and by noncompactness. The remaining
global-to-$S$-local kernel assembly in \S12.3, p.~514, must supply
a valid intertwiner and the required prime corrections; the
diagonal domination step alone no longer describes the gate. All-number-field Hecke scope and
zero-multiplicity assertions require additional proofs.
\item[Profinite GT, anabelian reconstruction and geometric spectra.]
The full fixed Ihara embedding's surjectivity, arithmetic associator
data, and a spectrum constructed geometrically from primes are distinct
targets. The current rational GT theorem and the adopted Dirichlet
theorem do not identify those objects. Explicit counterexamples here
refute only the stated routes through abelian cohomology or bare spectra.
\end{description}

\section{Checks and current status}
Every completed new proof is associated with its actual reading events
in \texttt{CHECK\_EVENTS.tsv}. No unproved candidate in this part is
used as an input to a completed proof. The timed checkpoints and the
three-hour, then hourly movement reviews are recorded in
\texttt{LEDGER.md}; the two-no-movement stopping rule remains in force.
